\guidechapter{Borrowed workspace and changing representations}{Read before the finite-network proof.}

Two ideas that initially seem unrelated work together here. First, a linear
calculation can borrow a register without knowing its contents, provided it
removes the influence of those contents and restores them. Second, arrays can
be represented in different coordinates at different places in a circuit.
The cost of changing representation can be smaller than directly performing
the desired overall array operation.

\paragraph{A four-update restoration example.}
Let $a$ be an arbitrary scratch value. Starting from $(x,y,a)$, perform
\[
 y\leftarrow y-a,\quad a\leftarrow a+x,\quad
 y\leftarrow y+a,\quad a\leftarrow a-x.
\]
The final state is $(x,y+x,a)$. The initial $a$ cancels from the target, and
the scratch register returns exactly to its initial value. This is a toy
identity over an additive ring, not a proposed implementation of the full
network. In characteristic two, addition and subtraction both become XOR.

\explfigure{03-restoration}{Unknown does not mean unusable.}{
Every intermediate occurrence of $a$ is shown. This small example illustrates
why starting scratch at zero would be an insufficient test.}

\paragraph{How the actual network isolates the right source.}
The original construction indexes values by triples $S,T\subseteq[h]$.
Gathering through coordinates shared by the triples gives a coefficient
depending on $|S\cap T|$. A side circuit cancels unwanted cross-terms. In
the bit version the central coefficient is $|S\cap T|$ modulo two:
intersection sizes zero and two already vanish, intersection one is canceled,
and the self-intersection three leaves one. Thus the net target update is
$y_S\leftarrow y_S+x_S$. For the complex version, the central coefficient
is $(|S\cap T|-1)/2$; side terms cancel the distinct intersection-zero and
intersection-two cases. In both cases the scratch-restoration schedule removes
the old auxiliary values.

Three bank updates yield
\[
 (x,y)\to(x,y+x)\to(-y,y+x)\to(-y,x).
\]
This is a signed swap; over bits the minus sign disappears. The network alone
has not yet accomplished a fast transform. It supplies an algebraic scaffold.

\paragraph{What a frame means.}
Let $g$ be the logical array on one role. A frame is an invertible linear
operator $D$ on its address-indexed entries. The stored array is $Dg$.
A pointwise gate, in contrast, combines different roles at the same address.
For a gate adding two roles, a common frame gives
\[
 Dg_1+Dg_2=D(g_1+g_2).
\]
This is ordinary linearity, but it requires the \emph{same} $D$ on every
incidence of that gate. With unrelated frames $D_1,D_2$, the desired identity
does not follow.

\explfigure{04-common-frame}{The diagram commutes because the frame is shared.}{
Horizontal arrows combine roles. Vertical arrows change the representation
within an array. These are different operations on different indices.}

Between vertices with frames $D_u,D_v$, the change is $D_vD_u^{-1}$.
Following the invariant through the circuit gives the implemented operator
\[
 D_{\rm out}\,\mathsf S_\Omega\,D_{\rm in}^{-1},
\]
where $\mathsf S_\Omega$ is the scalar circuit applied pointwise. If the scalar
circuit is a signed role permutation, each route sees only its endpoint frames
and its scalar sign. Intermediate frames cancel algebraically, but their
physical changes must still be paid for. Choosing them cheaply is the design
problem.

\paragraph{Why the geometry matters.}
The paper attaches subspaces to vertices. Comparable neighboring subspaces
give residual directions whose dimensions count smaller operations. Low rank
means fewer independent directions to implement. A matrix with many nonzero
entries may still have rank one; for example,
$\left(\begin{smallmatrix}1&1\\1&1\end{smallmatrix}\right)$ moves only
one independent combination. Conversely, a sparse identity matrix has full
rank. The accounting therefore depends on the actual residual operators,
their bases, and endpoint contributions, not just the number of visible gates.

The bit circuit combines scalars in $\mathbb F_2$ while using rational
subspace labels for address operators. The complex circuit combines Gaussian
dyadic scalars while using binary subspace labels for phase operators. These
are two different uses of a field; confusing scalar values with address labels
obscures the construction.

\takeaway{Restoration makes existing rows usable as workspace. Common frames
turn a scalar swap into an array operation. Residual ranks price the changes
between those frames. All three properties are needed.}
\sourceanchor{03-motifs.tex, lem:scalar-motifs; eq:common-frame-identity}
